\documentclass{article}
% Better font encoding for modern output
\usepackage{amssymb}
\usepackage[english]{babel}
\usepackage[T1]{fontenc}
\usepackage{lmodern}
\usepackage{blindtext}
\usepackage{layouts}
\usepackage{enumitem}

\title{\Large Lists}
\author{
	John Doe \\
	Department of Mathematics \\
	University of Science \\
	\texttt{john.doe@university.edu}
	\and
	Jane Smith \\
	Department of Physics \\
	University of Technology \\
	\texttt{jane.smith@university.edu}
}
\date{\today}

% Define custom text commands
\newcommand{\randomtext}{Hello, here is some text without a meaning. Hello, here is some text without a meaning. Hello, here is some text without a meaning.}
\newcommand{\shortrandomtext}{Short text without meaning.}

% Command with one argument
\newcommand{\highlight}[1]{\textbf{#1}}

% Command with default argument
\newcommand{\greeting}[1][Hello]{\textbf{#1}!}

% Command with multiple arguments
\newcommand{\person}[2]{#1 (#2 years old)}

% Define custom symbols for MetaPost replacements
\newcommand{\mpbullet}{$\bullet$}
\newcommand{\mpdot}{$\cdot$}
\newcommand{\mpyingyang}{$\bigcirc$}
\newcommand{\mpsquare}[3]{$\square$}

\begin{document}
	
	\maketitle
	
	PART I
	
	1. Lists
	Lists are easily to create:
	
	\begin{itemize}
		\item List entries start with the \verb|\item| command.
		\item Individual entries are indicated with a black dot, a so-called bullet.
		\item The text in the entries may be of any length.
	\end{itemize}
	
	2. Numbered (ordered) lists are easy to create:
	
	\begin{enumerate}
		\item Items are numbered automatically.
		\item The numbers start at 1 with each use of the \texttt{enumerate} environment.
		\item Another entry in the list
	\end{enumerate}
	
	3. Description Environment
	
	\begin{description}
		\item This is an entry \textit{without} a label.
		\item[Something short] A short one-line description.
		\item[Something long] A much longer description. \blindtext[1]
	\end{description}
	
	4. Labels
	Change the labels using \verb|\item[label text]| in an \texttt{itemize} environment
	
	\begin{itemize}
		\item This is my first point
		\item Another point I want to make 
		\item[!] A point to exclaim something!
		\item[$\blacksquare$] Make the point fair and square.
		\item[NOTE] This entry has no bullet
		\item[] A blank label?
	\end{itemize}
	
	\vspace{10pt}
	
	5. Nested Lists
	
	\begin{enumerate}
		\item The labels consists of sequential numbers
		\begin{itemize}
			\item The individual entries are indicated with a black dot, a so-called bullet
			\item The text in the entries may be of any length
			\begin{description}
				\item[Note:] I would like to describe something here
				\item[Caveat!] And give a warning here
			\end{description}
		\end{itemize}
		\item The numbers starts at 1 with each use of the \texttt{enumerate} environment
	\end{enumerate}
	
	6. Nested Lists: Label Style
	
	\begin{enumerate}
		\item First level item
		\item First level item
		\begin{enumerate}
			\item Second level item
			\item Second level item
			\begin{enumerate}
				\item Third level item
				\item Third level item
				\begin{enumerate}
					\item Fourth level item
					\item Fourth level item
				\end{enumerate}
			\end{enumerate}
		\end{enumerate}
	\end{enumerate}
	
	7. Nested itemize lists: bullet style
	
	\begin{itemize}
		\item First level item
		\item First level item
		\begin{itemize}
			\item Second level item
			\item Second level item
			\begin{itemize}
				\item Third level item
				\item Third level item
				\begin{itemize}
					\item Fourth level item
					\item Fourth level item
				\end{itemize}
			\end{itemize}
		\end{itemize}
	\end{itemize}
	
	PART II
	Customizing Lists 
	
	\renewcommand{\labelenumii}{\arabic{enumi}.\arabic{enumii}}
	\renewcommand{\labelenumiii}{\arabic{enumi}.\arabic{enumii}.\arabic{enumiii}}
	\renewcommand{\labelenumiv}{\arabic{enumi}.\arabic{enumii}.\arabic{enumiii}.\arabic{enumiv}}
	
	\begin{enumerate}
		\item One
		\item Two
		\item Three
		\begin{enumerate}
			\item Three point one
			\begin{enumerate}
				\item Three point one, point one
				\begin{enumerate}
					\item Three point one, point one, point one
					\item Three point one, point one, point two
				\end{enumerate}
			\end{enumerate}
		\end{enumerate}
		\item Four
		\item Five
	\end{enumerate}
	
	\begin{verbatim}
		\newcounter{foo}
		\setcounter{foo}{5}    
	\end{verbatim}
	\newcounter{foo}
	\setcounter{foo}{5} 
	
	\begin{itemize}
		\item \verb|\arabic{foo}| produces \arabic{foo}
		\item \verb|\roman{foo}| produces \roman{foo} 
		\item \verb|\Roman{foo}| produces \Roman{foo}
		\item \verb|\Alph{foo}| produces \Alph{foo}
		\item \verb|\alph{foo}| produces \alph{foo}
	\end{itemize}
	
	% Comment out the duck commands that don't exist
	% \renewcommand{\labelenumi}{\duck{enumi}}
	% \renewcommand{\labelenumii}{\duck{enumi}.\duckegg{enumii}}
	% \renewcommand{\labelenumiii}{\duck{enumi}.\duckegg{enumii}.\duckegg{enumiii}}
	% \renewcommand{\labelenumiv}{\duck{enumi}.\duckegg{enumii}.\duckegg{enumiii}.\duckchick{enumiv}} 
	
	\begin{enumerate}
		\item A duck
		\item More ducks
		\item A flurry of ducks
		\begin{enumerate}
			\item Ducks and eggs
			\begin{enumerate}
				\item Do I see... 
				\item Ducks and pre-ducks 
				\begin{enumerate}
					\item Awww...
					\item So cute!
				\end{enumerate}
			\end{enumerate}
		\end{enumerate}
		\item Back to ducks
		\item Again
	\end{enumerate}
	
	\begin{figure}
		\listdiagram  % This is from the layouts package - should work
		\caption{The \LaTeX{} parameters which define typesetting and layout of lists.} 
	\end{figure}
	
	\newcounter{boxlblcounter}  
	\newcommand{\makeboxlabel}[1]{\fbox{#1.}\hfill}% \hfill fills the label box
	\newenvironment{boxlabel}
	{\begin{list}
			{\arabic{boxlblcounter}}
			{\usecounter{boxlblcounter}
				\setlength{\labelwidth}{3em}
				\setlength{\labelsep}{0em}
				\setlength{\itemsep}{2pt}
				\setlength{\leftmargin}{1.5cm}
				\setlength{\rightmargin}{2cm}
				\setlength{\itemindent}{0em} 
				\let\makelabel=\makeboxlabel
			}
		}
		{\end{list}}
	
	\noindent\randomtext
	
	\begin{boxlabel}
		\item \randomtext
		\item \randomtext
		\item \randomtext
	\end{boxlabel}
	
	% Usage of custom commands
	\highlight{Important text}
	\greeting           % Outputs: Hello!
	\greeting[Hi]       % Outputs: Hi!
	\person{John}{25}   % Outputs: John (25 years old)
	
	\section*{Using LaTeX's default settings for \texttt{itemize}}
	
	\randomtext
	
	\begin{itemize}
		\item \randomtext
		\begin{itemize}
			\item \randomtext
			\begin{itemize}
				\item \randomtext
			\end{itemize}
			\item \randomtext
		\end{itemize}
		\item \randomtext
	\end{itemize}
	
	\section*{Using a custom \texttt{itemize} via \texttt{enumitem}}
	
	\subsection*{Note the effect of left and right margin settings}
	
	\randomtext
	
	\begin{itemize}[leftmargin=30pt, rightmargin=2cm]
		\item \randomtext
		\begin{itemize}
			\item \randomtext
			\begin{itemize}
				\item \randomtext
			\end{itemize}
			\item \randomtext
		\end{itemize}
		\item \randomtext
	\end{itemize}
	
	% Create a custom list based on enumerate
	% It is called "myitems"
	% We'll create a list that is 3 levels deep
	\newlist{myitems}{enumerate}{3}
	
	% Configure the behaviour of level 1 entries
	% NOTE: we use the list counter "myitemsi"
	\setlist[myitems, 1]
	{label=\arabic{myitemsi}., %1., 2., 3., ...
		leftmargin=\parindent,
		rightmargin=10pt
	}
	
	% Configure the behaviour of level 2 entries
	% NOTE: we use the list counter "myitemsii"
	\setlist[myitems, 2]
	{label=\arabic{myitemsi}.\arabic{myitemsii}, %1.1, 1.2, 1.3...
		leftmargin=15pt,
		rightmargin=15pt}
	
	% Configure the behaviour of level 3 entries
	% NOTE: we use the list counter "myitemsiii"
	\setlist[myitems, 3]
	% Use a label of 1.1:<kern>(a), 1.1:<kern>(b) etc  
	{label=\arabic{myitemsi}.\arabic{myitemsii}:\kern1.5pt(\alph{myitemsiii}),
		leftmargin=30pt,
		rightmargin=30pt}
	
	\randomtext
	
	% Now use the custom myitems environment
	\begin{myitems}
		\item \randomtext
		\begin{myitems}
			\item \randomtext
			\begin{myitems}
				\item \randomtext
				\item \randomtext
			\end{myitems}
			\item \shortrandomtext
			\item \shortrandomtext
		\end{myitems}
		\item \randomtext
	\end{myitems}
	
	\newlist{contract}{enumerate}{10}
	\setlist[contract]{label*=\arabic*.}
	\setlistdepth{10} 
	
	\section*{Custom list nested to 10 levels deep!}
	
	\begin{contract}
		\item Level 1
		\begin{contract}
			\item Level 2
			\begin{contract}
				\item Level 3
				\begin{contract}
					\item Level 4
					\begin{contract}
						\item Level 5
						\begin{contract}
							\item Level 6
							\begin{contract}
								\item Level 7
								\begin{contract}
									\item Level 8
									\begin{contract}
										\item Level 9
										\begin{contract}
											\item Level 10
										\end{contract}
									\end{contract}
								\end{contract}
							\end{contract}
						\end{contract}
					\end{contract}
				\end{contract}
			\end{contract}
		\end{contract}
	\end{contract}
	
	% Declare a new itemize-based list via enumitem
	\newlist{myEnumerate}{itemize}{6}
	
	% The nosep option removes all vertical spacing
	% the label=\protect\mpbullet causes all bullets to be 
	% drawn by a macro that uses MetaPost code. \protect
	% is required as noted in the enumitem manual 
	
	\setlist[myEnumerate]{nosep,label=\protect\mpbullet}
	\setlistdepth{6}
	
	\begin{myEnumerate}
		\item Level 1
		\begin{myEnumerate}
			\item Level 2
			\begin{myEnumerate}
				\item Level 3
				\begin{myEnumerate}
					\item Level 4
					\begin{myEnumerate}
						\item Level 5
						\begin{myEnumerate}
							\item Level 6
						\end{myEnumerate}
					\end{myEnumerate}
				\end{myEnumerate}
			\end{myEnumerate}
		\end{myEnumerate}
	\end{myEnumerate}
	
	% COMMENTED OUT - Requires MetaPost or additional packages
	% \begin{itemize}
		%     \item Start thinking about what we hope to achieve
		%     \begin{todolist}
			%         \item[\mpdot] Identify objectives
			%         \item[\mpyingyang] Balance environmental impact 
			%         \item[\mpsquare{0}{5}{0}] Implement plans
			%         \begin{todolist}
				%             \item[\mpsquare{-0.5}{4}{0}] Stage 1 plans
				%             \item[\mpsquare{-0.5}{4}{-20}] Stage 2 plans
				%             \item[\mpsquare{-0.5}{4}{-40}] Stage 3 plans
				%             \item[\mpsquare{-0.5}{4}{-60}] Stage 4 plans
				%         \end{todolist}
			%     \end{todolist}
		% \end{itemize}
	
	% ALTERNATIVE: Use standard symbols instead
	\begin{itemize}
		\item Start thinking about what we hope to achieve
		\begin{itemize}
			\item[$\bullet$] Identify objectives
			\item[$\circ$] Balance environmental impact 
			\item[$\square$] Implement plans
			\begin{itemize}
				\item[$\diamond$] Stage 1 plans
				\item[$\star$] Stage 2 plans
				\item[$\triangle$] Stage 3 plans
				\item[$\blacksquare$] Stage 4 plans
			\end{itemize}
		\end{itemize}
	\end{itemize}
	
\end{document}